\documentclass[11pt,a4paper]{article}
\usepackage[T1]{fontenc}\usepackage[utf8]{inputenc}\usepackage[english]{babel}
\usepackage{lmodern}\usepackage[a4paper,margin=2.5cm]{geometry}
\usepackage{amsmath,amssymb}\usepackage{siunitx}\usepackage{booktabs}
\usepackage{empheq}\usepackage{verbatim}
\usepackage[numbers,sort&compress]{natbib}
\usepackage[colorlinks=true,linkcolor=black,citecolor=black,urlcolor=blue]{hyperref}
\emergencystretch=1.5em

\begin{filecontents*}{pulsearea_refs.bib}
@book{goodman2017,
  author    = {Goodman, Joseph W.},
  title     = {Introduction to Fourier Optics},
  edition   = {4}, publisher = {W. H. Freeman / Macmillan Learning},
  address   = {New York}, year = {2017}, isbn = {978-1-319-11916-4},
  note      = {Fourier conventions, convolution and correlation theorems,
               Parseval's theorem; band-limited imaging}
}

@book{bornwolf1999,
  author    = {Born, Max and Wolf, Emil},
  title     = {Principles of Optics},
  edition   = {7 (expanded)}, publisher = {Cambridge University Press},
  address   = {Cambridge}, year = {1999}, isbn = {978-0-521-64222-4},
  note      = {Sec.~8.5, the Airy pattern of a circular aperture}
}

@misc{dlmf,
  author       = {{NIST}},
  title        = {{NIST} Digital Library of Mathematical Functions,
                  Chapter~10 (Bessel Functions), \S10.22 (Integrals)},
  howpublished = {\url{https://dlmf.nist.gov/10.22}},
  note         = {Release 1.2.4; accessed 17 September 2026}
}

@article{sibalic2017,
  author  = {\v{S}ibali\'{c}, N. and Pritchard, J. D. and Adams, C. S. and
             Weatherill, K. J.},
  title   = {{ARC}: An open-source library for calculating properties of
             alkali {R}ydberg atoms},
  journal = {Computer Physics Communications},
  volume  = {220}, pages = {319--331}, year = {2017},
  doi     = {10.1016/j.cpc.2017.06.015}
}

@article{grimm2000,
  author  = {Grimm, Rudolf and Weidem\"{u}ller, Matthias and
             Ovchinnikov, Yurii B.},
  title   = {Optical Dipole Traps for Neutral Atoms},
  journal = {Advances in Atomic, Molecular, and Optical Physics},
  volume  = {42}, pages = {95--170}, year = {2000},
  doi     = {10.1016/S1049-250X(08)60186-X},
  note    = {also \url{https://arxiv.org/abs/physics/9902072}}
}

@misc{steckrb85,
  author       = {Steck, Daniel A.},
  title        = {Rubidium 85 {D} Line Data},
  howpublished = {\url{https://steck.us/alkalidata}},
  note         = {accessed 17 September 2026}
}
\end{filecontents*}

\newcommand{\dd}{\mathrm d}
\title{Addendum: the correlation theorem behind the overlap integral,\\
       and sources for the pulse-area section}
\author{}\date{}
\begin{document}\maketitle

\section{Equation (42) in full}

The step in question is
\begin{equation}
O(\mathbf d)\;=\;\int\! u(|\mathbf r|)\,u(|\mathbf r-\mathbf d|)\,\dd^2r
\;=\;\frac{1}{(2\pi)^2}\int\!\bigl|\tilde u(q)\bigr|^2\,
      e^{\mathrm i\mathbf q\cdot\mathbf d}\,\dd^2q .
\label{eq:corr}
\end{equation}
It was quoted as "the correlation theorem''. Written out it is five lines, and
three of them are the ones that actually do work: the choice of Fourier
convention, the symmetry of $u$, and the interchange of the two integrations.

\paragraph{Convention.} The factor $1/(2\pi)^2$ in Eq.~\eqref{eq:corr} is not
universal -- it is fixed by where one puts it in the transform pair. Throughout,
\begin{equation}
\tilde u(\mathbf q)=\int\! u(\mathbf r)\,e^{-\mathrm i\mathbf q\cdot\mathbf r}\,\dd^2r ,
\qquad
u(\mathbf r)=\frac{1}{(2\pi)^2}\int\!\tilde u(\mathbf q)\,
             e^{+\mathrm i\mathbf q\cdot\mathbf r}\,\dd^2q ,
\label{eq:conv}
\end{equation}
i.e.\ the whole $1/(2\pi)^2$ sits in the back transform \cite[Ch.~2]{goodman2017}.
This is the same convention used for $\tilde u(q)=(4\pi/k^2)\,\Theta(k-q)$ in
step~(i), so the two are consistent.

\paragraph{Step 1: replace the shifted factor by its transform.}
Only the second factor is expanded; the first is left alone.
\begin{equation}
u(\mathbf r-\mathbf d)
=\frac{1}{(2\pi)^2}\int\!\tilde u(\mathbf q)\,
 e^{\mathrm i\mathbf q\cdot(\mathbf r-\mathbf d)}\,\dd^2q
=\frac{1}{(2\pi)^2}\int\!\tilde u(\mathbf q)\,
 e^{\mathrm i\mathbf q\cdot\mathbf r}\,e^{-\mathrm i\mathbf q\cdot\mathbf d}\,\dd^2q .
\end{equation}
The shift has become a pure phase factor $e^{-\mathrm i\mathbf q\cdot\mathbf d}$
-- this is the only property of the Fourier transform the derivation needs.

\paragraph{Step 2: insert and exchange the order of integration.}
\begin{align}
O(\mathbf d)
&=\int\!\dd^2r\; u(\mathbf r)\,\frac{1}{(2\pi)^2}
  \int\!\dd^2q\;\tilde u(\mathbf q)\,e^{\mathrm i\mathbf q\cdot\mathbf r}\,
  e^{-\mathrm i\mathbf q\cdot\mathbf d}\\[2pt]
&=\frac{1}{(2\pi)^2}\int\!\dd^2q\;\tilde u(\mathbf q)\,
  e^{-\mathrm i\mathbf q\cdot\mathbf d}\;
  \underbrace{\int\!\dd^2r\; u(\mathbf r)\,e^{+\mathrm i\mathbf q\cdot\mathbf r}}
            _{\textstyle =\;\tilde u(-\mathbf q)} .
\label{eq:step2}
\end{align}
The inner integral is the defining transform \eqref{eq:conv} evaluated at
$-\mathbf q$, because the sign in the exponent is the opposite one.
Interchanging the two integrations is legitimate here: $u$ is bounded and
square integrable and $\tilde u$ has compact support, so the double integral
converges absolutely and Fubini's theorem applies. (For a Gaussian mode the
same statement is immediate; for the Airy mode it is the compact support of
$\tilde u$ that saves it, since $u$ itself decays only as $r^{-3/2}$.)

\paragraph{Step 3: the symmetry of the mode.} Two properties are needed, and
the mode has both:
\begin{align}
u \ \text{real} \quad&\Longrightarrow\quad
  \tilde u(-\mathbf q)=\tilde u^{*}(\mathbf q) &&\text{(Hermitian symmetry),}\\
u \ \text{even},\ u(-\mathbf r)=u(\mathbf r) \quad&\Longrightarrow\quad
  \tilde u(-\mathbf q)=\tilde u(\mathbf q) &&\text{(substitute } \mathbf r\to-\mathbf r).
\end{align}
A radially symmetric mode is automatically even, and both the Gaussian and
$2J_1(kr)/(kr)$ are real. Taken together, $\tilde u$ is itself real and even --
as the explicit result $\tilde u=(4\pi/k^2)\Theta(k-q)$ shows. Therefore
\begin{equation}
\tilde u(\mathbf q)\,\tilde u(-\mathbf q)
=\tilde u(\mathbf q)\,\tilde u^{*}(\mathbf q)
=\bigl|\tilde u(\mathbf q)\bigr|^{2} ,
\end{equation}
and Eq.~\eqref{eq:step2} becomes
\begin{equation}
O(\mathbf d)=\frac{1}{(2\pi)^2}\int\!\bigl|\tilde u(\mathbf q)\bigr|^{2}\,
             e^{-\mathrm i\mathbf q\cdot\mathbf d}\,\dd^2q .
\label{eq:step3}
\end{equation}

\paragraph{Step 4: the sign in the exponent does not matter.}
Equation~\eqref{eq:step3} carries $e^{-\mathrm i\mathbf q\cdot\mathbf d}$ while
Eq.~\eqref{eq:corr} was written with $e^{+\mathrm i\mathbf q\cdot\mathbf d}$.
Both are right: $|\tilde u|^2$ is even in $\mathbf q$, so the substitution
$\mathbf q\to-\mathbf q$ leaves the measure and the weight unchanged and only
flips the exponent. Equivalently: $O$ is the transform of a real even function
and is therefore itself real and even, $O(-\mathbf d)=O(\mathbf d)$ -- which is
also physically obvious, since the overlap of two spots cannot depend on which
one is called the first.

\paragraph{Why "correlation'' and "convolution'' coincide here.}
In general the two differ,
\begin{equation}
(u\star v)(\mathbf d)=\int\! u(\mathbf r)\,v(\mathbf r-\mathbf d)\,\dd^2r ,
\qquad
(u * v)(\mathbf d)=\int\! u(\mathbf r)\,v(\mathbf d-\mathbf r)\,\dd^2r ,
\end{equation}
and only the convolution turns into a plain product of transforms. For an
\emph{even} $v$ one has $v(\mathbf d-\mathbf r)=v(\mathbf r-\mathbf d)$ and the
two integrals are identical \cite[Ch.~2]{goodman2017}. That is the one line the
phrase "since $u$ is real and even the correlation equals the convolution''
stands for.

\paragraph{Step 5: the check at \boldmath$\mathbf d=0$.}
Setting $\mathbf d=0$ in Eq.~\eqref{eq:corr} reduces it to Parseval's theorem
in this convention,
\begin{equation}
O(0)=\int\! u^2\,\dd^2r
    =\frac{1}{(2\pi)^2}\int\!|\tilde u|^2\,\dd^2q \;=\;P_1 ,
\end{equation}
so the normalisation of Eq.~\eqref{eq:corr} is fixed by a quantity already
computed independently in step~(a). This is the cheapest available test of the
prefactor and it passes.

\paragraph{Evaluation for the Airy mode.}
Now the special structure enters. An indicator function is idempotent,
$\Theta^2=\Theta$, and that is what makes the Airy case as easy as it is. With
$\tilde u(q)=(4\pi/k^2)\,\Theta(k-q)$,
\begin{align}
O^{\mathrm A}(d)
&=\frac{1}{(2\pi)^2}\left(\frac{4\pi}{k^2}\right)^{\!2}
  \int_{|\mathbf q|<k}\! e^{\mathrm i\mathbf q\cdot\mathbf d}\,\dd^2q
&&\text{Eq.~\eqref{eq:step3}, }\Theta^2=\Theta\\
&=\frac{1}{4\pi^{2}}\cdot\frac{16\pi^{2}}{k^{4}}
  \cdot\frac{2\pi k\,J_1(kd)}{d}
&&\text{disc integral, step (ii)}\\
&=\frac{4}{k^{4}}\cdot\frac{2\pi k\,J_1(kd)}{d}
 \;=\;\frac{8\pi\,J_1(kd)}{k^{3}d}\\
&=\underbrace{\frac{4\pi}{k^{2}}}_{P_1^{\mathrm A}}
  \cdot\underbrace{\frac{2J_1(kd)}{kd}}_{u(d)} .
\end{align}
\begin{empheq}[box=\fbox]{equation}
  O^{\mathrm A}(d)=P_1^{\mathrm A}\,u(d)
\end{empheq}
The autocorrelation of the Airy mode is the mode itself. The reason is visible
in the first line and nowhere else: the spectrum is flat inside the cutoff, so
squaring it does nothing, and a function whose spectrum equals that of $u$ up
to a constant \emph{is} $u$ up to that constant. No such statement holds for
the Gaussian, whose squared spectrum is narrower by $\sqrt2$ -- which is
exactly why $O^{\mathrm G}(d)=P_1^{\mathrm G}e^{-d^2/2w_0^2}$ carries
$2w_0^2$ in the exponent instead of $w_0^2$.

\paragraph{Numerical confirmation.}
On a $2048^2$ grid over \SI{40}{\micro\meter} with $w_0=\SI{1.04}{\micro\meter}$:
$O(0)$ reproduces $P_1$ to machine precision (as it must, both being the same
sum), the analytic $P_1^{\mathrm A}=4\pi/k^2$ is met to \SI{1.2}{\percent}, and
$O(d)$ agrees with $P_1^{\mathrm A}u(d)$ to \SI{1.1}{\percent} of the peak over
$d = 0.2\dots\SI{6}{\micro\meter}$ -- the residual being the grid truncating the
outer Airy rings, consistent with the \SI{5}{\percent} truncation quoted in the
main text. At the working-point separation $d=\SI{3.305}{\micro\meter}$ one gets
$u(d)=+0.06308$.

\section{Suggested citations for the pulse-area section}

None of the derivations need a source; the starting points do. In the order
they appear:

\begin{center}
\begin{tabular}{p{0.44\textwidth}p{0.48\textwidth}}
\toprule
Place in the text & Citation\\
\midrule
"$\int_0^\infty J_\nu^2(x)\,\dd x/x = 1/(2\nu)$'' (Airy spot power) &
\verb|\cite{dlmf}| -- Weber--Schafheitlin type, DLMF \S10.22\\[4pt]
"$\int_0^\infty J_1(x)\,\dd x = 1$'' (prefactor of $\tilde u$) &
\verb|\cite{dlmf}|\\[4pt]
"$\int_0^X J_0(t)\,t\,\dd t = X J_1(X)$'' (disc integral) &
\verb|\cite{dlmf}|; equivalently one line from
$\frac{\dd}{\dd z}\!\left[zJ_1(z)\right]=zJ_0(z)$\\[4pt]
"the focal field is the Fourier transform of the pupil field'' and the
Airy form $2J_1(kr)/(kr)$ &
\verb|\cite[Ch.~4,~6]{goodman2017}|, \verb|\cite[Sec.~8.5]{bornwolf1999}|\\[4pt]
the correlation theorem, Eq.~(42), and Parseval &
\verb|\cite[Ch.~2]{goodman2017}|\\[4pt]
"the excited states are eliminated adiabatically'' and the two-photon
Rabi frequency &
\verb|\cite{grimm2000}| for the light shift and scattering rate in the same
far-detuned limit\\[4pt]
signed dipole matrix elements, hyperfine shifts, line centroids, lifetimes &
\verb|\cite{sibalic2017}| (ARC), cross-checked against \verb|\cite{steckrb85}|\\[4pt]
the thermal width
$\sigma^2=\frac{\hbar}{2m\omega_r}\coth\frac{\hbar\omega_r}{2k_BT}$ &
\verb|\cite{tuchendler2008}| (already in the appendix bibliography)\\
\bottomrule
\end{tabular}
\end{center}

\noindent Two remarks on what \emph{not} to cite. The boxed pulse-area formula
and the cancellation of $A_\mathrm{eff}$ against $\tilde P_\mathrm{prof}$ are
bookkeeping carried out here in full -- a reference would suggest the result is
someone else's. The same holds for the sinc: it is the Fourier transform of the
integration window and is derived in three lines in the text.

\bibliographystyle{unsrtnat}
\bibliography{pulsearea_refs}

\clearpage
\section*{BibTeX source (additional entries)}
\noindent These are new; \texttt{tuchendler2008} is already in the appendix
bibliography.
\scriptsize
\verbatiminput{pulsearea_refs.bib}
\normalsize

\end{document}
